\section{Support-relative dependency}
\label{sec:support-dependency}

The continuous example in \cref{sec:supported-process} already shows that coordinate dependence is not determined by an ambient formula alone. This section records the exact criterion for coordinate sufficiency and one consequence for circuit-style necessity tests.

For $K\subseteq I_q$, let
\[
X_{q,K}:=\prod_{i\in K}A_{q,i}
\]
and let $\pi_K:X_q\to X_{q,K}$ be coordinate projection. For $K=\varnothing$, take $X_{q,\varnothing}$ to be a singleton.

\begin{definition}[Sufficient coordinate set]
A subset $K\subseteq I_q$ is \emph{sufficient for receiver $j$ on $C$} if there exists
\[
\widehat Q_{q,j,K}^{\theta,C}:
\pi_K(S_q^{\theta,C})\to Z_{q,j}^{\theta,C}
\]
such that
\begin{equation}
Q_{q,j}^{\theta,C}
=
\widehat Q_{q,j,K}^{\theta,C}
\,\pi_K|_{S_q^{\theta,C}}.
\label{eq:coordinate-sufficiency}
\end{equation}
\end{definition}

The definition asks whether the receiver state can be reconstructed from the coordinates in $K$ on the supported domain. It does not ask whether the ambient expression for $Q_{q,j}^\theta$ contains those coordinates syntactically. Relations that hold only on the support can make one coordinate determine another.

\begin{proposition}[Coordinate sufficiency criterion]
\label{prop:coordinate-sufficiency}
A coordinate set $K$ is sufficient for receiver $j$ on $C$ if and only if
\begin{equation}
\Ker(\pi_K|_{S_q^{\theta,C}})
\subseteq
\Ker Q_{q,j}^{\theta,C}.
\label{eq:coordinate-kernel-criterion}
\end{equation}
\end{proposition}

\begin{proof}
If \eqref{eq:coordinate-sufficiency} holds, equal projected values produce equal receiver values. Conversely, suppose \eqref{eq:coordinate-kernel-criterion} holds. Define
\[
\widehat Q_{q,j,K}^{\theta,C}(\pi_K(x))
:=Q_{q,j}^\theta(x).
\]
The kernel inclusion makes this independent of the representative $x\in S_q^{\theta,C}$.
\end{proof}

Because $I_q$ is finite and $I_q$ itself is sufficient, every receiver has at least one inclusion-minimal sufficient coordinate set. Write
\[
\MinSupp_{q,j}(\theta,C)
\]
for the collection of these minimal sets.

Minimality inherits the domain dependence of sufficiency. On the diagonal in the continuous example, either coordinate is a minimal sufficient access for the linear readout even though both are required on $\mathbb R^2$. Restriction can also remove a dependency entirely. For
\[
Q(a,b)=a\oplus b
\]
on $\{0,1\}^2$, both coordinates are needed on the full square, while $Q$ is constant on the diagonal $\{(0,0),(1,1)\}$. The empty coordinate set is therefore sufficient after restriction.

More generally, sufficiency is monotone under restriction. If $S\subseteq S'\subseteq X_q$ and $K$ is sufficient for a fixed map $Q$ on $S'$, then $K$ is sufficient on $S$. The same kernel inclusion is simply being tested on fewer pairs. Minimality need not be preserved, because a proper subset of $K$ may become sufficient on the smaller domain.

\subsection{One-coordinate operativity and necessity}

A common local test varies one component while holding the others fixed and asks whether a downstream quantity changes. On a restricted support, such a comparison is available only when both states belong to the support.

\begin{definition}[Direct supported operativity]
For coordinate $i\in I_q$ and receiver $j$, write
\[
i\to_{\theta,C}j
\]
when there exist $x,x'\in S_q^{\theta,C}$ such that
\[
x_k=x'_k\quad\text{for all }k\neq i,
\qquad
Q_{q,j}^\theta(x)\neq Q_{q,j}^\theta(x').
\]
\end{definition}

\begin{proposition}[Direct operativity and minimal sufficient supports]
\label{prop:direct-necessity}
For fixed $(q,j,\theta,C)$ and $i\in I_q$,
\begin{equation}
i\to_{\theta,C}j
\quad\Longleftrightarrow\quad
i\in K
\text{ for every }
K\in\MinSupp_{q,j}(\theta,C).
\label{eq:direct-necessity}
\end{equation}
\end{proposition}

\begin{proof}
Suppose $i\to_{\theta,C}j$, witnessed by $x,x'$. Any coordinate set $K$ omitting $i$ gives $\pi_K(x)=\pi_K(x')$ while the receiver values differ, so $K$ is not sufficient. Every minimal sufficient set must therefore contain $i$.

Conversely, suppose no direct witness exists. Whenever two supported states agree on all coordinates except possibly $i$, their receiver values agree. Hence $I_q\setminus\{i\}$ is sufficient by \cref{prop:coordinate-sufficiency}. Since $I_q$ is finite, it contains an inclusion-minimal sufficient subset omitting $i$.
\end{proof}

The proposition separates several circuit claims that are easy to conflate. A one-coordinate witness characterizes membership in every inclusion-minimal sufficient coordinate set. Membership in at least one minimal set, attribution thresholds, and interventions performed on a larger counterfactual domain are different criteria.
